\section{The Bipartisan Case}
\label{sec:bipartisan}

\subsection{Why Both Parties Have Reason to Accept the DIA}

Federal redistricting mandates have historically failed not for lack of a
mathematical solution but for lack of a political coalition.
Republicans resist mandates that they perceive as designed to reduce their
seat count; Democrats resist mandates that would reduce theirs.
Every proposed algorithm has been attacked as partisan by whichever party
it disfavors in the current cycle.

The \dia{} is designed to be genuinely acceptable to both parties---not
by splitting the difference between their preferences, but by satisfying
each party's \emph{legitimate} interest in a redistricting system that
cannot be weaponized against them.

The argument to each party is different because each party's legitimate
interest is different.
We address them separately.

\subsection{The Republican Case}

Republicans' core objection to federal redistricting mandates has been, in
order of importance: (1) racial targeting that packs non-Hispanic white voters
to dilute their influence, (2) algorithms that produce non-compact,
bizarrely shaped districts in urban areas, (3) mandates that lock in
Democratic advantages in high-population states, and (4) federal
overreach into a traditionally state function.

The \dia{} addresses each concern directly.

\paragraph{No racial targeting.}
The default \ar{} run contains zero racial data.
The inputs are: the prime factorization of $k$, the TIGER adjacency graph,
and census tract populations.
The TIGER graph is a geometric description of which tracts share a boundary.
The population data is a count of persons in each tract.
Neither input contains any racial identification.

The algorithm minimizes edge cut---total shared boundary length between
districts.
This criterion produces compact districts.
It does not, by construction or by outcome, target any racial group.

This is not an evasion.
\callais{} (2026) establishes the legal framework for algorithmic
redistricting, and it is precisely the framework that \ar{} satisfies.
The majority at page 36 requires that ``race-conscious and partisan signals
must be disentangled.''
\ar{}'s default run uses \emph{neither} signal---it is below the threshold
of either racial-conscious districting or partisan-conscious districting.
A court applying the \shaw{}--\miller{} racial-predominance test to an
\ar{} map has nothing to find, because race was not an input.

\paragraph{Compact districts.}
The minimum-edge-cut criterion is a mathematically rigorous formulation of
compactness.
It minimizes the total length of district boundaries, which is equivalent
to minimizing the perimeter-to-area ratio across the map as a whole.
The resulting districts are, in the language of \textit{Polsby and Popper}
\citep{polsby1991}, compact: their boundaries follow natural geographic
contours rather than the population-concentration boundaries that
characterize racial or partisan gerrymanders.

The edge-weighted version of the algorithm (using shared boundary lengths
as edge weights, from \citet{b2paper}) produces an additional 56\%
improvement in the Polsby-Popper compactness score over the unweighted
version.
Republican objections to oddly shaped districts---the objections raised
against majority-minority districts and against Democratic urban gerrymanders
alike---are addressed by this criterion.

\paragraph{County preservation by default.}
The \dia{} specifies a default county-preservation weight of $\alpha_\text{county} = 2.0$,
implemented via the B.10 subdivision-respecting edge-weighting scheme
\citep{b10paper}.
At this weight, county boundaries are treated as strong natural lines:
the algorithm strongly prefers to place an entire county in a single
district rather than splitting it across district boundaries.
County-line preservation is a traditionally Republican concern, reflecting
the role of counties as the basic unit of Republican local political
organization.
The statutory default encodes this preference while allowing states that
wish to use a different weight to petition for a variance through the
Election Assistance Commission.

\paragraph{No locked-in Democratic advantage.}
The bipartisan concern about locked-in advantages cuts both ways.
Republicans who currently control map-drawing in their states worry that
a federal mandate will reduce their advantage.
But the advantage they currently hold was purchased through a process that
Democrats will eventually reverse when they gain legislative control in those
states.
The cycle is costly, litigious, and produces maps that serve no one's
constituents.

\ar{} breaks the cycle by removing the control that either party can exercise.
In a state like Wisconsin, where the current Republican map has been found
unconstitutional by the state Supreme Court, \ar{} produces a map that
neither party drew and neither party can attack as partisan.
The Wisconsin bakeoff shows that \ar{} produces a 4D/4R split---which, in a
state with a roughly 50/50 partisan environment, is the outcome that most
closely tracks the state's actual partisan composition.

More importantly: a statute mandating \ar{} \emph{cannot be used later to
favor Democrats}, because the algorithm has no partisan inputs.
When Democrats next hold legislative majorities in red-leaning states, they
cannot use \ar{} to gerrymander---not because the statute forbids it, but
because the algorithm's inputs (factorization of $k$ and TIGER graph) do not
contain the partisan data required for partisan manipulation.
The protection is structural, not merely statutory.

\paragraph{Federal authority has a well-established constitutional basis.}
Article~I, \S4 of the Constitution provides that ``The Times, Places and
Manner of holding Elections for Senators and Representatives, shall be
prescribed in each State by the Legislature thereof; but the Congress may
at any time by Law make or alter such Regulations.''\footnote{%
  U.S. Const., Art.\ I, \usca{}~4, cl.\ 1.
}
The Elections Clause gives Congress plenary authority to regulate the
manner of congressional elections, including the redistricting process.
The Supreme Court confirmed this in \textit{Smiley v.\ Holm},
285 U.S. 355, 366--67 (1932), and again in \textit{Arizona State
Legislature v.\ Arizona Independent Redistricting Commission},
576 U.S. 787 (2015).\footnote{%
  \textit{Arizona State Legislature v.\ Arizona Independent Redistricting
  Commission}, 576 U.S. 787, 813 (2015): ``Congress has the power to
  regulate congressional elections, and the Elections Clause gives Congress
  authority that covers redistricting.''
}
A Republican objection to federal authority over redistricting would have
to run against 85 years of settled constitutional law.

\subsection{The Democratic Case}

Democrats' core objections to federal redistricting mandates have been:
(1) algorithms that systematically undercount urban and minority populations,
(2) VRA non-compliance that dilutes minority voting strength, (3) mandates
that lock in Republican structural advantages in rural-majority states,
and (4) census-stability requirements that prevent courts from correcting
constitutional violations mid-decade.

The \dia{} addresses each concern.

\paragraph{Strict population equality.}
The $\pm 0.05\%$ population tolerance (\texttt{ubvec[0]=1.001}) is an
order of magnitude tighter than any redistricting plan that has survived
judicial review.
\karcher{} tolerance is effectively zero, but courts have upheld plans
with deviations of less than 1\%.
The \dia{} goes further: its maximum district-population deviation is
smaller than the Census Bureau's margin of error in counting individual tracts.
Democratic concerns about minority undercounting are addressed at the
input stage, not the algorithm stage: if the Census Bureau's decennial count
undercounts minority communities, the remedy is census methodology reform,
not redistricting algorithm manipulation.
But the algorithm will apply the correct population equality standard to
whatever population data the Census Bureau provides.

\paragraph{VRA compliance via the Gingles overlay.}
The \dia{} mandates a separate VRA overlay process for states where the
Gingles three-prong test is met \citep{b14paper}.
The overlay, implemented via VRASection, uses geographic alignment scores
to prefer bisection ratios that allow geographically compact minority
communities to form majority-minority districts.

This is not racial targeting in the \shaw{}--\miller{} sense.
The alignment score measures where minority population is geographically
concentrated, not how many minority voters the district should contain.
It is a geographic property of the existing population distribution.
The Gingles prong-1 requirement---that the minority group must be
``sufficiently large and geographically compact'' to constitute a majority
in a single-member district---is a geographic test.
VRASection operationalizes it as a geographic measurement.

\allenmilligan{} (2023) required Alabama to draw two Black-opportunity
districts.
The VRASection overlay, calibrated to Alabama's 2020 census data, produces
exactly this outcome: two districts with Black voting-age population
above 50\%, confirmed empirically by \citet{b14paper}.\footnote{%
  \citet{b14paper}: Alabama 2:5 result confirmed.
  Two majority-minority districts at 55.2\% and 52.8\% Black VAP.
  VRA weight $w_\text{vra} = 0.40$ calibrated below the \shaw{}--\miller{}
  racial-predominance threshold.
}

The VRA overlay is Callais-compliant because it contains no partisan
signal: VRASection uses only the geographic alignment of minority VAP
concentration, not any measurement of how minority voters vote.
The \callais{} disentanglement requirement is satisfied at the architectural
level.

\paragraph{No structural Republican advantage.}
The concern that algorithmic redistricting systematically advantages
Republicans reflects the empirical reality that geographic voter sorting
\citep{rodden2019,chen2013}---the concentration of Democratic voters in
urban areas---creates a Republican structural advantage under any compact-district
algorithm.
This is the Rodden-Chen result: compact districts, drawn without political
inputs, tend to produce Republican-leaning outcomes in states with
significant urban Democratic concentrations.

This concern is real, but it does not argue against the \dia{}.
It argues for proportional representation, which requires a change to the
single-member-district system that is itself a statutory default, not a
constitutional requirement.
The \dia{} proposes the best redistricting algorithm within the
single-member-district system.
If Democrats wish to challenge the single-member-district system, the
appropriate vehicle is different legislation.

More practically: \ar{}'s Wisconsin outcome (4D/4R) is better for Democrats
than GeoSection's outcome (3D/5R) in that state, as the bakeoff confirms.
The algorithm that is constitutionally derived is also, in Wisconsin, the
one that produces a more proportional outcome.
This is not coincidental: the \hh{} bisection tree divides the state into
geographic regions based on the natural factorization of the seat count,
which in Wisconsin (8 = 2$\times$2$\times$2) produces a four-quadrant split
that roughly tracks the state's political geography.
The constitutional algorithm happens to be more proportional than the
practitioner's algorithm in the most-litigated state in the country.

\paragraph{Census stability does not prevent constitutional correction.}
The \dia{}'s census-stability provision (\citet{b15paper}: 67\% of states
produce maps that remain valid across all three census years) does not prevent
courts from ordering corrective redistricting when a constitutional violation
occurs.
The statute provides that maps shall remain in effect for the decade, but
this provision is subject to the courts' remedial authority.
A court finding a constitutional violation may order a new \ar{} run using
updated data.
The statutory stability provision addresses voluntary mid-decade redrawing,
not court-ordered remediation.

\paragraph{The process is transparent and auditable.}
Every decision in the \ar{} process is documented in the \texttt{whatif-manifest
v1} JSON sidecar produced by the \texttt{redist} CLI \citep{b0paper}.
The manifest records every parameter used, the seed (derivable from the
census release identifier), the TIGER graph version, and the resulting
district boundaries.
A civil rights organization, a Democratic advocacy group, or a court
can independently reproduce the map from the census data and the
manifest in approximately 30 minutes using the open-source \texttt{redist}
binary.
There is no black box.
There is no discretion.
There is no room for the algorithm to be run in a way that produces a
different output.
The audit trail is complete and independently verifiable.

\subsection{The Courts Will Accept It}

Courts have struggled with redistricting because the legal standards for
evaluating maps are contested.
Partisan gerrymandering is non-justiciable at the federal level after \rucho{}.
Racial gerrymandering is justiciable but the standard---racial predominance
under \shaw{}--\miller{}---requires a factual inquiry into the mapmaker's
intent that is difficult to resolve on the record.
Compactness is a criterion in most state statutes but courts have declined
to make it a constitutional floor.

The \dia{} sidesteps all of these difficulties.

\paragraph{No new constitutional doctrine required.}
The first-principles argument in Section~\ref{sec:first-principles} extends
an existing constitutional standard---the Huntington-Hill method mandated by
2~U.S.C.~\usca{}~2a(a)---to the next level of the apportionment hierarchy.
Courts do not need to invent new standards.
They need only to recognize that the same principle applies at the intrastate
level as at the interstate level.
This is the kind of doctrinal extension that courts find most comfortable:
applying an established rule to a new factual context.

\paragraph{The racial-predominance test is moot.}
When the mapmaker is an algorithm with no racial inputs, the question
``did race predominate?'' has an obvious answer: no.
The \shaw{}--\miller{} standard asks whether ``race was the predominant factor
overriding traditional districting principles.''
In an \ar{} map, the only inputs are population, adjacency, and the seat-count
factorization.
Race is not an input.
The racial-predominance inquiry is foreclosed by the algorithm's design,
not by the court's deference.

\paragraph{Partisan gerrymandering is moot.}
\rucho{} held that partisan gerrymandering is non-justiciable as a federal
constitutional matter.
\ar{} maps contain no partisan inputs.
There is no partisan gerrymandering to adjudicate.
The entire class of Rucho-type litigation disappears.

\paragraph{The content-derived seed forecloses seed-manipulation claims.}
One potential avenue for challenging any algorithmic map is that the
mapmaker chose a seed that produced a favorable outcome.
This concern is addressed by the statutory seed formula:
\texttt{SHA-256(census\_release\_id || `DIA\_SEED\_V1')}.
The seed is computed from the Census Bureau's public release identifier
for the decennial census data, which is known at the moment the census
data is released.
No one---not the mapmaker, not the controlling party, not the Census Bureau---
can retroactively choose a seed to favor a particular outcome.
The seed is content-derived, publicly computable, and verifiable.
This forecloses any claim that the seed was manipulated.

\paragraph{The strong-inference test under Callais is satisfied.}
\callais{} at page 23 establishes that \S2 challengers must show that the
state's stated political goals could have been achieved with better minority
outcomes.
The \dia{} explicitly provides this counterfactual capacity.
A challenger can run the VRASection overlay---the \callais{}-compliant
strong-inference alternative---on the same data and show that better minority
outcomes are achievable.
The toolbox provides both the default map and the VRA-optimized alternative
in a form that courts can directly compare.
The \callais{} strong-inference framework is not just satisfied; it is operationalized.
